% Russian supplement, embedded-inclusion version; values in guillemets are filled by fill_embedded.py.
\documentclass[11pt,a4paper]{article}
\usepackage[T2A]{fontenc}
\usepackage[utf8]{inputenc}
\usepackage[english,russian]{babel}
\usepackage[margin=25mm]{geometry}
\usepackage{amsmath,amssymb,graphicx}
\usepackage{placeins}
\usepackage[numbers,sort&compress]{natbib}
\usepackage[hyphens]{url}
\usepackage[hidelinks]{hyperref}
\setlength{\emergencystretch}{2em}
\ifx\pdfmapfile\undefined\else
\pdfmapfile{+cm-super-t1.map}\pdfmapfile{+cm-super-t2a.map}
\pdfmapfile{+cm-super-ts1.map}\fi
\renewcommand{\thesection}{S\arabic{section}}
\renewcommand{\theequation}{S\arabic{equation}}
\renewcommand{\thefigure}{S\arabic{figure}}
\renewcommand{\thetable}{S\arabic{table}}

\title{Дополнительные материалы к статье <<Упругий отклик алюминия на магнитострикцию включения Al$_{13}$Fe$_4$
и размер зоны возможной пластичности>>}
\author{И.\,Е.~Михайловский, Д.\,Е.~Пшонкин}
\date{}

\begin{document}
\maketitle

\section{Построение и подготовка ячейки}

Матрица строится как периодический кристалл алюминия с параметром решётки
$a_0=4{,}05$~\AA{}, при котором объёмный алюминий в этом потенциале
не напряжён. Вдоль $x\parallel[1\bar10]$ ячейка содержит 44 периода
по 2{,}864~\AA{}, вдоль $y\parallel[11\bar2]$ 25 периодов по 4{,}960~\AA{},
вдоль $z\parallel[111]$ 54 слоя $(111)$ по 2{,}338~\AA{}; размеры ячейки
$126{,}0\times124{,}0\times126{,}3$~\AA{}. Включение вырезано шаром радиусом
25~\AA{} из кристалла Al$_{13}$Fe$_4$, ось $a$ которого направлена вдоль $x$,
а ось $b$ вдоль $y$. Во включении 4\,487 атомов, из них 3\,429 Al и 1\,058 Fe
(отношение 3{,}24 при стехиометрическом 3{,}25). Из матрицы удалены атомы
алюминия внутри шара и все атомы ближе 2{,}2~\AA{} к атомам включения
(4\,426 атомов). Наименьшие межатомные расстояния после построения:
2{,}20~\AA{} между матрицей и включением, 2{,}37~\AA{} внутри включения.

С этим потенциалом кристалл Al$_{13}$Fe$_4$ при свободной релаксации
меняет решётку по сравнению с экспериментальной: $a$ на $+7{,}5\%$,
$b$ на $-2{,}5\%$, высоту ячейки на $+3{,}0\%$, объём на $+7{,}9\%$.
Поэтому на всех этапах атомы включения закреплены в экспериментальной
решётке. В пробном расчёте со свободным включением ячейка при 300~К
вытягивалась по осям неравномерно, что указывает на раздувание включения
внутри алюминия; этот расчёт был остановлен и далее не используется.

Подготовка отсчётной ячейки: минимизация энергии матрицы до евклидовой нормы
сил 1~эВ/\AA{}; 10~пс при 300~К (термостат Нозе--Гувера, только атомы
матрицы, размер ячейки постоянный); охлаждение до 0{,}5~К за 10~пс
(термостат Ланжевена); минимизация порциями по 2000 итераций до нормы сил
$3\cdot10^{-4}$~эВ/\AA{}. После этого по атомам матрицы дальше 15~\AA{}
от поверхности включения вычисляется среднее нормальное напряжение
$\sigma_0=«CTL_FAR_MEAN»$~МПа. Все размеры ячейки умножаются на
$1-\sigma_0/3K=«CTL_BOX_SCALE»$, где модуль всестороннего сжатия алюминия
в потенциале $K=77{,}4$~ГПа; при этом смещаются только атомы матрицы.
Затем минимизация повторяется, и среднее напряжение вдали от включения
становится равным $«CTL_FAR_AFTER»$~МПа.

\section{Сходимость и линейность}

«S2_TEXT»

\section{Заданная деформация и свободная магнитострикция}

Форма включения в расчёте задана, поэтому деформация
(1) основной статьи и есть фактическая деформация включения внутри матрицы.
Поле в матрице определяется только смещением поверхности включения
и упругостью матрицы. Пусть упругая сферическая частица имеет модуль сдвига
$k\mu$, где $\mu$ обозначает модуль сдвига алюминия, и свободную
магнитострикцию $\lambda$ той же формы, что и (1), то есть деформацию,
которую частица приняла бы в поле без матрицы. Её фактическая деформация
внутри матрицы меньше $\lambda$. Решение Эшелби для неоднородного включения,
упругие постоянные которого отличаются от постоянных матрицы
\cite{Eshelby1957,Mura1987}, при девиаторной деформации (не меняющей объёма,
как (1)) даёт
\begin{equation}
\varepsilon=A\lambda,\qquad A=\frac{\beta k}{1+\beta(k-1)},\qquad
\beta=\frac{2(4-5\nu)}{15(1-\nu)} ,
\end{equation}
где коэффициент Пуассона $\nu$ для простоты принят одинаковым для включения
и матрицы; при $\nu=0{,}344$ получается $\beta=0{,}463$. Если $\lambda$ задана,
все напряжения в матрице равны рассчитанным, умноженным на $A\lambda/\varepsilon$.
Для включения той же жёсткости, что и матрица ($k=1$), $A=0{,}46$; для
включения в 2 и 2{,}5 раза жёстче $A=0{,}63$ и $0{,}68$; для абсолютно
жёсткого включения $A=1$. Упругие постоянные Al$_{13}$Fe$_4$ из этого
потенциала рассчитаны для равновесной в нём решётки Al$_{13}$Fe$_4$,
которая завышена (раздел~S1), и для оценки $k$ не используются.

\section{Решения теории упругости}

\paragraph{Одиночная сфера.} Сферическое включение радиуса $R$ с однородной
деформацией превращения $\boldsymbol\varepsilon^*$ (деформацией, которую
включение приняло бы, если бы его вынули из матрицы) в бесконечной изотропной
матрице имеет внутри однородную деформацию
$\boldsymbol\varepsilon=\mathbf S:\boldsymbol\varepsilon^*$, где
$\mathbf S$ обозначает тензор Эшелби \cite{Eshelby1957,Mura1987}. Чтобы внутренняя
деформация совпала с заданной (1), берётся
$\boldsymbol\varepsilon^*=\mathbf S^{-1}:\boldsymbol\varepsilon$;
для девиаторной деформации это $\boldsymbol\varepsilon/\beta$. Деформация
снаружи, $\varepsilon_{ij}(\mathbf x)=D_{ijkl}(\mathbf x)\,\varepsilon^*_{kl}$,
выражается через внешний тензор Эшелби \cite{JuSun1999}; при $\rho=R/r$
и $\mathbf n=\mathbf x/r$
\begin{multline}
30(1-\nu)\rho^{-3}D_{ijkl}=(3\rho^2+10\nu-5)\delta_{ij}\delta_{kl}
+(3\rho^2-10\nu+5)(\delta_{ik}\delta_{jl}+\delta_{il}\delta_{jk})
+15(1-\rho^2)\delta_{ij}n_kn_l\\
+15(1-2\nu-\rho^2)\delta_{kl}n_in_j
+15(\nu-\rho^2)(\delta_{ik}n_jn_l+\delta_{jl}n_in_k+\delta_{il}n_jn_k+\delta_{jk}n_in_l)
+15(7\rho^2-5)n_in_jn_kn_l .
\end{multline}
Напряжение находится по закону Гука
$\sigma_{ij}=\lambda_L\varepsilon_{kk}\delta_{ij}+2\mu\varepsilon_{ij}$
с изотропными постоянными алюминия из потенциала ($\mu=26{,}9$~ГПа,
$\nu=0{,}344$, усреднение Фойгта--Ройсса--Хилла по $C_{11}$, $C_{12}$,
$C_{44}$). Программа проверена на задаче о центре дилатации
($\boldsymbol\varepsilon^*=10^{-3}\mathbf I$): при $r=1{,}5R$ она даёт
$\sigma_{rr}=-21{,}77$ и $\sigma_{\theta\theta}=10{,}89$~МПа, что совпадает
с точным решением Ляве.

\paragraph{Та же периодическая ячейка.} Уравнения равновесия
для периодической решётки включений с периодом ячейки решаются в пространстве
Фурье \cite{MoulinecSuquet1998}. Для однородной среды с жёсткостью
$\mathbf C$ (кубические постоянные алюминия из потенциала, повёрнутые
к осям ячейки) и поля деформации превращения $\chi(\mathbf x)
\boldsymbol\varepsilon^*$, где $\chi$ равна единице во включении,
\begin{equation}
\hat\varepsilon_{ij}(\mathbf k)=\tfrac12\bigl[\xi_iN_{jm}(\boldsymbol\xi)
+\xi_jN_{im}(\boldsymbol\xi)\bigr]\tau_{ml}\xi_l\,\hat\chi(\mathbf k),
\qquad \mathbf N=(\boldsymbol\xi\cdot\mathbf C\cdot\boldsymbol\xi)^{-1},
\quad \boldsymbol\tau=\mathbf C:\boldsymbol\varepsilon^*,
\end{equation}
$\boldsymbol\xi=\mathbf k/|\mathbf k|$; член $\mathbf k=0$ опущен, то есть
средняя деформация ячейки равна нулю, как в расчёте с неизменным размером.
Сетка с шагом 1~\AA{}, граница включения сглажена на одну ячейку сетки.
Деформация $\boldsymbol\varepsilon^*$ подбирается так, чтобы средняя
деформация внутри включения равнялась заданной (1): в периодической ячейке
для этого её нужно увеличить в 1{,}07--1{,}11 раза (в зависимости
от компоненты) по сравнению с одиночной сферой.
Тот же расчёт даёт смещения, которые используются как начальное
приближение для атомов матрицы в ячейке $+\varepsilon$. Среднее
дополнительное напряжение в матрице при неизменном размере ячейки известно
точно, $\langle\Delta\boldsymbol\sigma\rangle=-\mathbf C:\boldsymbol\varepsilon\,
f/(1-f)$; при $f=3{,}3\%$ это $«MEAN_PRED»$~МПа.

\section{Модель ползучести}

Дислокация преодолевает препятствие благодаря тепловым колебаниям атомов;
частота таких событий равна $\omega_0\exp(-\Delta G/k_{\mathrm B}T)$, где
$\omega_0$ обозначает частоту попыток, а $\Delta G$ высоту энергетического
барьера. Касательное
напряжение $\tau$ совершает на пути через барьер работу $V^*\tau$, поэтому
барьер для прыжков вперёд понижается, а для прыжков назад повышается:
$\Delta G_\pm=\Delta G_0\mp V^*\tau$ \cite{KocksArgonAshby1975,CaillardMartin2003}.
Если в объёме вокруг включений встречаются напряжения обоих знаков
с равными весами, средний множитель скорости равен
$\tfrac12[\exp(x)+\exp(-x)]=\cosh x$, $x=V^*\tau/k_{\mathrm B}T$.
Когда включения разрежены и их поля не перекрываются, относительный
прирост средней скорости в материале с объёмной долей включений $f$ равен
интегралу $\cosh x-1$ по объёму матрицы вокруг одного включения,
умноженному на число включений в единице объёма $f/V_p$; это формула (2) основной статьи. Для поля вида $\tau(r)=\tau_m(R/r)^3$, где $\tau_m$ обозначает касательное
напряжение у поверхности включения, замена
$u=x_0(R/r)^3$ сводит интеграл к
\begin{equation}
H=f\,x_0\int_0^{x_0}\frac{\cosh u-1}{u^2}\,\mathrm du,\qquad
x_0=\frac{V^*\tau_m}{k_{\mathrm B}T},
\end{equation}
по этой формуле построена сплошная кривая рис.~6. Штриховая кривая
получена из рассчитанного поля одиночной сферы: интеграл (2) взят
по сферической сетке (120 радиальных узлов, равномерно распределённых по $(R/r)^3$, что
охватывает всю матрицу, и $48\times48$ узлов по углам) с касательными напряжениями на двенадцати системах скольжения
$\{111\}\langle110\rangle$ для 64 случайных ориентаций зерна; затем
найдена амплитуда $\tau_m$, при которой формула (S4) даёт тот же $H$.
Прирост на 25\% достигается при $V^*=«VSTAR_GRAINS»\,b^3$ (усреднение
по ориентациям зёрен), «VSTAR_CRYSTAL»$\,b^3$ (одна ориентация, как
в расчётной ячейке) и «VSTAR_MAX»$\,b^3$ (для той же ориентации в каждой
точке берётся наибольшее из касательных напряжений на двенадцати системах
скольжения).
Поскольку поле пропорционально $\varepsilon$, $H$ зависит только
от произведения $V^*\varepsilon$, и при $V^*=70\,b^3$ прирост 25\%
даёт деформация $\varepsilon\cdot«VSTAR_GRAINS»/70\approx«EPS_AT_70»\%$.

\section{Сопротивление решётки движению дислокации}

Чтобы оценить, можно ли атомистически проследить движение дислокации мимо
включения, измерено, при каком напряжении в этом потенциале скользит
прямая краевая дислокация без включения. Ячейка: пластина алюминия
$126{,}0\times19{,}8\times125{,}5$~\AA{} (18\,768 атомов), периодическая
вдоль $x$ и $y$, со свободными поверхностями по $z$; по три верхних и нижних
слоя $(111)$ жёстко закреплены. Краевая дислокация с вектором Бюргерса
$b=2{,}864$~\AA{} вдоль $x$ и линией вдоль $y$ введена в плоскость
скольжения на высоте 91{,}2~\AA{}: в верхней половине пластины на одну
колонку атомов вдоль $x$ больше, чем в нижней. Верхний блок сдвигается
вдоль $x$ шагами по 0{,}00455~\AA, после каждого шага энергия минимизируется
до нормы сил $10^{-3}$~эВ/\AA{} при 0~К. Приложенное касательное напряжение
равно силе, с которой матрица действует на верхний и нижний блоки,
делённой на площадь $L_xL_y$ (среднее по двум блокам, они различаются не
больше чем на 1~МПа).

Напряжение растёт линейно, пока дислокация стоит, и падает, когда она
делает шаг в соседнее положение с наименьшей энергией. Длина шага
1{,}41--1{,}42~\AA, то есть $b/2$; положения линии, найденные
по смещению блока, совпадают с найденными методом DXA \cite{Stukowski2012DXA}
в пределах 0{,}25~\AA, и DXA во всех кадрах находит две частичные
дислокации: в алюминии краевая дислокация расщепляется на две, разделённые
полосой, где нарушен порядок укладки плоскостей $(111)$. Перед каждым из трёх перескоков напряжение доходит до 9{,}13,
9{,}00 и 9{,}01~МПа. Таким образом, в этом потенциале прямая краевая
дислокация при 0~К начинает скользить при $\tau_P\approx9$~МПа (напряжение
Пайерлса).

«S6_TEXT»

\section{Данные и программы}
\sloppy

Программы и входные файлы (пути от корня репозитория):
\begin{itemize}
\item построение ячейки: \path{scripts/build_embedded_inclusion.py};
деформированные ячейки: \path{scripts/make_strained_inclusion.py};
входные файлы LAMMPS: \path{lammps/embedded_inclusion/in.relax_control}
и \path{lammps/embedded_inclusion/in.static_pair};
\item напряжения: \path{analysis/python/embedded_inclusion_field.py};
проверки сходимости и линейности: \path{analysis/python/embedded_gates.py};
поиск дефектов: \path{analysis/python/embedded_dxa.py}; деформация,
напряжение на единицу деформации и плотность у границы:
\path{analysis/python/embedded_interface.py};
решения теории упругости: \path{analysis/python/embedded_inclusion_continuum.py};
\item модель ползучести: \path{analysis/python/embedded_creep.py};
порог текучести у частицы: \path{analysis/python/embedded_yield.py};
упругие постоянные потенциала: \path{analysis/python/meam_elastic_constants.py};
рисунки: \path{analysis/python/embedded_figures.py};
\item пластина с дислокацией: \path{scripts/build_embedded_dislocation.py};
релаксация: \path{lammps/embedded_inclusion/in.disl_relax};
положения дислокации: \path{analysis/python/embedded_disl_states.py};
сдвиг: \path{lammps/embedded_inclusion/in.disl_sweep}
и \path{analysis/python/embedded_sweep.py};
\item запуск расчётов ячеек $\pm\varepsilon$: \path{scripts/queue_embedded_static.sh};
расчёты после них (дислокация над включением, ячейки с деформацией
$4\varepsilon$ и $10\varepsilon$, сдвиг): \path{scripts/queue_embedded_plasticity.sh};
записи расчётов с дислокацией и ячеек $4\varepsilon$ и $10\varepsilon$:
\path{analysis/python/embedded_records.py}.
\end{itemize}
Записи результатов: \path{docs/reports/embedded_field.json},
\path{docs/reports/embedded_gates.json}, \path{docs/reports/embedded_continuum.json},
\path{docs/reports/embedded_creep.json}, \path{docs/reports/embedded_yield.json},
\path{docs/reports/embedded_peierls.json}, \path{docs/reports/embedded_dxa.json},
\path{docs/reports/embedded_interface.json}, \path{docs/reports/embedded_margin.json},
\path{docs/reports/embedded_disl_states.json},
\path{docs/reports/meam_elastic_constants.json}. Конфигурации и журналы
расчётов в сжатом виде с контрольными суммами лежат
в \path{data/publication/embedded}.

\bibliographystyle{unsrtnat}
\bibliography{references}
\end{document}
